% Section IV (part) -- Implementation: baseline models (Person 2)
% Include inside Section IV of main.tex with:  % Section IV (part) -- Implementation: baseline models (Person 2)
% Include inside Section IV of main.tex with:  % Section IV (part) -- Implementation: baseline models (Person 2)
% Include inside Section IV of main.tex with:  % Section IV (part) -- Implementation: baseline models (Person 2)
% Include inside Section IV of main.tex with:  \input{sections/baselines}
% Parameter counts below are for DEFAULT_CONFIG in config.py; re-check them with
%   python -m tests.test_baselines
% if any hyperparameter changes.
\subsection{Baseline Models}
\label{sec:baselines}

To isolate the effect of the attention mechanism, all three models share
the same convolutional front end, positional encoding, encoder-block wrapper,
pooling and classification head (\texttt{models/common.py}); they differ
\emph{only} in the sequence-mixing block. Every model maps an input window
$X \in \mathbb{R}^{B \times T \times 40}$ to logits in
$\mathbb{R}^{B \times 3}$ and accepts any window length $T$ without
re-instantiation.

\textbf{Shared CNN backbone.} Two 1D convolutional layers (kernel size 3,
``same'' padding, 64 channels) operate along the time axis, each followed
by batch normalisation, GELU and dropout ($p=0.1$). Because padding
preserves $T$, every downstream block sees the full history length.

\textbf{Classification head.} Features are mean-pooled over time, then
passed through LayerNorm, dropout and a linear layer to three logits
(Up, Stationary, Down).

\textbf{Baseline 1: CNN.} The backbone is connected directly to the head.
With $d_{\text{model}}=64$ this model has 20{,}675 trainable parameters.

\textbf{Baseline 2: CNN + standard attention.} After the backbone we add
fixed sinusoidal positional encodings and one pre-norm Transformer encoder
block: $h \leftarrow h + \mathrm{Drop}(\mathrm{Attn}(\mathrm{LN}(h)))$,
$h \leftarrow h + \mathrm{Drop}(\mathrm{FFN}(\mathrm{LN}(h)))$, with a
two-layer GELU feed-forward network of width $2 d_{\text{model}}$.
Attention uses $H=4$ heads of size $d_h = 16$:
\begin{equation}
  \mathrm{Attn}(Q,K,V) = \mathrm{softmax}\!\left(\frac{QK^{\top}}{\sqrt{d_h}}\right) V ,
\end{equation}
which materialises a $T \times T$ score matrix per head, giving $O(T^2 d)$
time and $O(H T^2)$ activation memory. The model has 54{,}147 parameters;
the encoder block therefore adds 33{,}472 parameters on top of Baseline~1.
The proposed model uses an identical block with linear attention
substituted (Section~\ref{sec:proposed}).

\textbf{Implementation note.} PyTorch's fused
\texttt{scaled\_dot\_product\_attention} may dispatch to memory-efficient
kernels that never store the $T \times T$ matrix. To keep the efficiency
comparison like-for-like with our plain-PyTorch linear attention, all
reported runs use the explicit implementation above; a unit test confirms
that both paths produce numerically identical outputs ($\le 10^{-5}$).

\textbf{Training.} All models are trained with the shared configuration:
AdamW (learning rate $10^{-3}$, weight decay $10^{-4}$), batch size 64,
up to 30 epochs, cross-entropy loss with inverse-frequency class weights
computed on the training split only, and early stopping (patience 5) on
validation macro-F1. The checkpoint with the best validation macro-F1 is
evaluated once on the held-out test split. Random seeds for Python, NumPy
and PyTorch are fixed before model and data-loader construction.

% Parameter counts below are for DEFAULT_CONFIG in config.py; re-check them with
%   python -m tests.test_baselines
% if any hyperparameter changes.
\subsection{Baseline Models}
\label{sec:baselines}

To isolate the effect of the attention mechanism, all three models share
the same convolutional front end, positional encoding, encoder-block wrapper,
pooling and classification head (\texttt{models/common.py}); they differ
\emph{only} in the sequence-mixing block. Every model maps an input window
$X \in \mathbb{R}^{B \times T \times 40}$ to logits in
$\mathbb{R}^{B \times 3}$ and accepts any window length $T$ without
re-instantiation.

\textbf{Shared CNN backbone.} Two 1D convolutional layers (kernel size 3,
``same'' padding, 64 channels) operate along the time axis, each followed
by batch normalisation, GELU and dropout ($p=0.1$). Because padding
preserves $T$, every downstream block sees the full history length.

\textbf{Classification head.} Features are mean-pooled over time, then
passed through LayerNorm, dropout and a linear layer to three logits
(Up, Stationary, Down).

\textbf{Baseline 1: CNN.} The backbone is connected directly to the head.
With $d_{\text{model}}=64$ this model has 20{,}675 trainable parameters.

\textbf{Baseline 2: CNN + standard attention.} After the backbone we add
fixed sinusoidal positional encodings and one pre-norm Transformer encoder
block: $h \leftarrow h + \mathrm{Drop}(\mathrm{Attn}(\mathrm{LN}(h)))$,
$h \leftarrow h + \mathrm{Drop}(\mathrm{FFN}(\mathrm{LN}(h)))$, with a
two-layer GELU feed-forward network of width $2 d_{\text{model}}$.
Attention uses $H=4$ heads of size $d_h = 16$:
\begin{equation}
  \mathrm{Attn}(Q,K,V) = \mathrm{softmax}\!\left(\frac{QK^{\top}}{\sqrt{d_h}}\right) V ,
\end{equation}
which materialises a $T \times T$ score matrix per head, giving $O(T^2 d)$
time and $O(H T^2)$ activation memory. The model has 54{,}147 parameters;
the encoder block therefore adds 33{,}472 parameters on top of Baseline~1.
The proposed model uses an identical block with linear attention
substituted (Section~\ref{sec:proposed}).

\textbf{Implementation note.} PyTorch's fused
\texttt{scaled\_dot\_product\_attention} may dispatch to memory-efficient
kernels that never store the $T \times T$ matrix. To keep the efficiency
comparison like-for-like with our plain-PyTorch linear attention, all
reported runs use the explicit implementation above; a unit test confirms
that both paths produce numerically identical outputs ($\le 10^{-5}$).

\textbf{Training.} All models are trained with the shared configuration:
AdamW (learning rate $10^{-3}$, weight decay $10^{-4}$), batch size 64,
up to 30 epochs, cross-entropy loss with inverse-frequency class weights
computed on the training split only, and early stopping (patience 5) on
validation macro-F1. The checkpoint with the best validation macro-F1 is
evaluated once on the held-out test split. Random seeds for Python, NumPy
and PyTorch are fixed before model and data-loader construction.

% Parameter counts below are for DEFAULT_CONFIG in config.py; re-check them with
%   python -m tests.test_baselines
% if any hyperparameter changes.
\subsection{Baseline Models}
\label{sec:baselines}

To isolate the effect of the attention mechanism, all three models share
the same convolutional front end, positional encoding, encoder-block wrapper,
pooling and classification head (\texttt{models/common.py}); they differ
\emph{only} in the sequence-mixing block. Every model maps an input window
$X \in \mathbb{R}^{B \times T \times 40}$ to logits in
$\mathbb{R}^{B \times 3}$ and accepts any window length $T$ without
re-instantiation.

\textbf{Shared CNN backbone.} Two 1D convolutional layers (kernel size 3,
``same'' padding, 64 channels) operate along the time axis, each followed
by batch normalisation, GELU and dropout ($p=0.1$). Because padding
preserves $T$, every downstream block sees the full history length.

\textbf{Classification head.} Features are mean-pooled over time, then
passed through LayerNorm, dropout and a linear layer to three logits
(Up, Stationary, Down).

\textbf{Baseline 1: CNN.} The backbone is connected directly to the head.
With $d_{\text{model}}=64$ this model has 20{,}675 trainable parameters.

\textbf{Baseline 2: CNN + standard attention.} After the backbone we add
fixed sinusoidal positional encodings and one pre-norm Transformer encoder
block: $h \leftarrow h + \mathrm{Drop}(\mathrm{Attn}(\mathrm{LN}(h)))$,
$h \leftarrow h + \mathrm{Drop}(\mathrm{FFN}(\mathrm{LN}(h)))$, with a
two-layer GELU feed-forward network of width $2 d_{\text{model}}$.
Attention uses $H=4$ heads of size $d_h = 16$:
\begin{equation}
  \mathrm{Attn}(Q,K,V) = \mathrm{softmax}\!\left(\frac{QK^{\top}}{\sqrt{d_h}}\right) V ,
\end{equation}
which materialises a $T \times T$ score matrix per head, giving $O(T^2 d)$
time and $O(H T^2)$ activation memory. The model has 54{,}147 parameters;
the encoder block therefore adds 33{,}472 parameters on top of Baseline~1.
The proposed model uses an identical block with linear attention
substituted (Section~\ref{sec:proposed}).

\textbf{Implementation note.} PyTorch's fused
\texttt{scaled\_dot\_product\_attention} may dispatch to memory-efficient
kernels that never store the $T \times T$ matrix. To keep the efficiency
comparison like-for-like with our plain-PyTorch linear attention, all
reported runs use the explicit implementation above; a unit test confirms
that both paths produce numerically identical outputs ($\le 10^{-5}$).

\textbf{Training.} All models are trained with the shared configuration:
AdamW (learning rate $10^{-3}$, weight decay $10^{-4}$), batch size 64,
up to 30 epochs, cross-entropy loss with inverse-frequency class weights
computed on the training split only, and early stopping (patience 5) on
validation macro-F1. The checkpoint with the best validation macro-F1 is
evaluated once on the held-out test split. Random seeds for Python, NumPy
and PyTorch are fixed before model and data-loader construction.

% Parameter counts below are for DEFAULT_CONFIG in config.py; re-check them with
%   python -m tests.test_baselines
% if any hyperparameter changes.
\subsection{Baseline Models}
\label{sec:baselines}

To isolate the effect of the attention mechanism, all three models share
the same convolutional front end, positional encoding, encoder-block wrapper,
pooling and classification head (\texttt{models/common.py}); they differ
\emph{only} in the sequence-mixing block. Every model maps an input window
$X \in \mathbb{R}^{B \times T \times 40}$ to logits in
$\mathbb{R}^{B \times 3}$ and accepts any window length $T$ without
re-instantiation.

\textbf{Shared CNN backbone.} Two 1D convolutional layers (kernel size 3,
``same'' padding, 64 channels) operate along the time axis, each followed
by batch normalisation, GELU and dropout ($p=0.1$). Because padding
preserves $T$, every downstream block sees the full history length.

\textbf{Classification head.} Features are mean-pooled over time, then
passed through LayerNorm, dropout and a linear layer to three logits
(Up, Stationary, Down).

\textbf{Baseline 1: CNN.} The backbone is connected directly to the head.
With $d_{\text{model}}=64$ this model has 20{,}675 trainable parameters.

\textbf{Baseline 2: CNN + standard attention.} After the backbone we add
fixed sinusoidal positional encodings and one pre-norm Transformer encoder
block: $h \leftarrow h + \mathrm{Drop}(\mathrm{Attn}(\mathrm{LN}(h)))$,
$h \leftarrow h + \mathrm{Drop}(\mathrm{FFN}(\mathrm{LN}(h)))$, with a
two-layer GELU feed-forward network of width $2 d_{\text{model}}$.
Attention uses $H=4$ heads of size $d_h = 16$:
\begin{equation}
  \mathrm{Attn}(Q,K,V) = \mathrm{softmax}\!\left(\frac{QK^{\top}}{\sqrt{d_h}}\right) V ,
\end{equation}
which materialises a $T \times T$ score matrix per head, giving $O(T^2 d)$
time and $O(H T^2)$ activation memory. The model has 54{,}147 parameters;
the encoder block therefore adds 33{,}472 parameters on top of Baseline~1.
The proposed model uses an identical block with linear attention
substituted (Section~\ref{sec:proposed}).

\textbf{Implementation note.} PyTorch's fused
\texttt{scaled\_dot\_product\_attention} may dispatch to memory-efficient
kernels that never store the $T \times T$ matrix. To keep the efficiency
comparison like-for-like with our plain-PyTorch linear attention, all
reported runs use the explicit implementation above; a unit test confirms
that both paths produce numerically identical outputs ($\le 10^{-5}$).

\textbf{Training.} All models are trained with the shared configuration:
AdamW (learning rate $10^{-3}$, weight decay $10^{-4}$), batch size 64,
up to 30 epochs, cross-entropy loss with inverse-frequency class weights
computed on the training split only, and early stopping (patience 5) on
validation macro-F1. The checkpoint with the best validation macro-F1 is
evaluated once on the held-out test split. Random seeds for Python, NumPy
and PyTorch are fixed before model and data-loader construction.
